%-------------------------
% Based off of: https://github.com/sb2nov/resume
% License : MIT
%------------------------

\documentclass[letterpaper,10pt]{article}

\usepackage{latexsym}
\usepackage[empty]{fullpage}
\usepackage{titlesec}
\usepackage{marvosym}
\usepackage[usenames,dvipsnames]{color}
\usepackage{verbatim}
\usepackage{enumitem}
\usepackage[hidelinks]{hyperref}
\usepackage{fancyhdr}
\usepackage[russian, english]{babel}
\usepackage{tabularx}
\input{glyphtounicode}


%----------FONT OPTIONS----------
% sans-serif
% \usepackage[sfdefault]{FiraSans}
% \usepackage[sfdefault]{roboto}
% \usepackage[sfdefault]{noto-sans}
% \usepackage[default]{sourcesanspro}

% serif
% \usepackage{CormorantGaramond}
% \usepackage{charter}


\pagestyle{fancy}
\fancyhf{} % clear all header and footer fields
\fancyfoot{}
\renewcommand{\headrulewidth}{0pt}
\renewcommand{\footrulewidth}{0pt}

% Adjust margins
\addtolength{\oddsidemargin}{-0.5in}
\addtolength{\evensidemargin}{-0.5in}
\addtolength{\textwidth}{1in}
\addtolength{\topmargin}{-.5in}
\addtolength{\textheight}{1.0in}

\urlstyle{same}

\raggedbottom
\raggedright
\setlength{\tabcolsep}{0in}

% Sections formatting
\titleformat{\section}{
	\vspace{-4pt}\scshape\raggedright\large
}{}{0em}{}[\color{black}\titlerule \vspace{-5pt}]

% Ensure that generate pdf is machine readable/ATS parsable
\pdfgentounicode=1

%-------------------------
% Custom commands
\newcommand{\resumeItem}[2]{
	\item\small{
		{\ifx&#1&\else\textbf{#1}: \fi#2 \vspace{-2pt}}
	}
}

\usepackage[T2A]{fontenc}
\usepackage[utf8]{inputenc}
\usepackage{cmap}

\newcommand{\resumeSubheading}[4]{
	\vspace{-1pt}\item
	\begin{tabular*}{0.97\textwidth}{l@{\extracolsep{\fill}}r}
		\textbf{#1} & #2 \\
		\ifx&#3&\else\textit{\small#3} & \textit{\small #4} \\ \fi
	\end{tabular*}\vspace{-5pt}
}


\newcommand{\resumeSubSubheading}[2]{
	\item
	\begin{tabular*}{0.97\textwidth}{l@{\extracolsep{\fill}}r}
		\textit{\small#1} & \textit{\small #2} \\
	\end{tabular*}\vspace{-7pt}
}

\newcommand{\resumeProjectHeading}[2]{
	\item
	\begin{tabular*}{0.97\textwidth}{l@{\extracolsep{\fill}}r}
		\small#1 & #2 \\
	\end{tabular*}\vspace{-7pt}
}

\newcommand{\resumeSubItem}[2]{\resumeItem{#1}{#2}\vspace{-4pt}}

\renewcommand\labelitemii{$\vcenter{\hbox{\tiny$\bullet$}}$}

\newcommand{\resumeSubHeadingListStart}{\begin{itemize}[leftmargin=0.15in, label={}]}
	\newcommand{\resumeSubHeadingListEnd}{\end{itemize}}
\newcommand{\resumeItemListStart}{\begin{itemize}}
	\newcommand{\resumeItemListEnd}{\end{itemize}\vspace{-5pt}}
\newcommand{\ulhref}[2]{\href{#1}{\underline{#2}}}
%-------------------------------------------
%%%%%%  RESUME STARTS HERE  %%%%%%%%%%%%%%%%%%%%%%%%%%%%


\begin{document}
	
	%----------HEADING----------
	% \begin{tabular*}{\textwidth}{l@{\extracolsep{\fill}}r}
		%   \textbf{\href{http://sourabhbajaj.com/}{\Large Sourabh Bajaj}} & Email : \href{mailto:sourabh@sourabhbajaj.com}{sourabh@sourabhbajaj.com}\\
		%   \href{http://sourabhbajaj.com/}{http://www.sourabhbajaj.com} & Mobile : +1-123-456-7890 \\
		% \end{tabular*}
	
	\begin{center}
		\textbf{\Large \scshape Усатов Павел} \textbf{\Large $|$ Java Backend Developer} \\ \vspace{1pt}
		\small +7 (912) 712-83-52 $|$ 
		\href{https://t.me/Usatov_Pavel}{\underline{Telegram}} $|$
		\href{mailto:pvusatov@edu.hse.ru}{\underline{pvusatov@edu.hse.ru}} $|$
		\href{https://github.com/UsatovPavel}{\underline{Github}} $|$
        \href{https://usatovpavel.ru}{\underline{\textbf{Сайт}}}
	\end{center}
	
	
	%-----------EDUCATION-----------
	\section{Образование}
	\resumeSubHeadingListStart
	\resumeSubheading
	{ВШЭ СПБ}{Санкт-Петербург}
	{Прикладная математика и информатика}{2023-2027}
	\item \textbf{Релевантные курсы:}
	Архитектура компьютера и операционные системы, Java, Scala, Kotlin, C++, базы данных,
	алгоритмы и структуры данных, компьютерные сети, функциональное программирование, архитектура ПО
	%технологии промышленной разработки ПО
	\resumeSubHeadingListEnd
	
	
	%-----------EXPERIENCE-----------
	\section{Опыт}
	\resumeSubHeadingListStart
	\resumeSubheading
	{Стажировка Backend Т-Банк}{апрель 2026 - июль 2026}
	{Участвовал в миграции документов с long на uuid}{Продуктовая команда документооборота}
	\resumeItemListStart
	\resumeItem{}{Версионировал 10 эндпоинтов под uuid, чтобы не сломать мобильные приложения, установленные у пользователей}
	\resumeItem{}{Написал модуль команды в общем репозитории интеграции b2b: 7 клиентов с документацией взамен legacy трёхлетней давности. Смежные команды завели эпики перехода на новое API}
	\resumeItem{}{Перевёл 23 интеграции low-code на новое API, обеспечив переход бухгалтерии. Проверил e2e на QA-контуре}
	\resumeItem{}{Работа с БД: создание алертов, оптимизация запроса, добавление расширений}
	\resumeItem{}{Чинил баги генерации документов: время у нас, восстановил столбец статистики во внешнем сервисе}
    \resumeItemListEnd
	\textbf{Технологии}: Scala, Cats Effect, PostgreSQL, Kanban, Kafka, Docker, Typescript, ELK
	\resumeSubheading
	{\ulhref{https://github.com/UsatovPavel/riid}{Утилита RIID для внутреннего облака VK}}{январь 2026 - июнь 2026}
	{Daemon на Java для p2p-загрузки OCI/Docker образов \textit{независимый} от container engine}{Курсовая, научрук - SRE из VK}
	\resumeItemListStart
    \resumeItem{}{{Написал \ulhref{https://github.com/UsatovPavel/java-dragonfly-image-puller}{библиотеку}: gRPC-клиент к Dragonfly для загрузки слоёв с соседних нод. \textbf{Взят в использование в VK}}}
    \resumeItem{}{Обосновал команде интеграцию  с Rust версией вместо согласованной ранее Go: Go не поддерживается с 2026}
	\resumeItem{}{Модульная архитектура: CLI $\rightarrow$ dispatcher (стратегия скачивания) $\rightarrow$ registry client/cache/p2p $\rightarrow$ engine adapters}
	\resumeItem{}{На 20\% быстрее Podman на кластере k8s из 12 нод: 100 образов от 1 МБ до 5 ГБ}
	%\resumeItem{}{Сбор метрик через Prometheus, отображение в Grafana}
	\resumeItemListEnd
	\textbf{Технологии}: Java, Kubernetes, Grafana, Prometheus, Gradle, OCI/Docker Registry API, p2p, Podman, Docker, gRPC%\ \ \ \ \ \ \ \ \ \ \ \ \ \ \ \ HTTP, CLI, интеграционные и e2e тесты
	
	
	\resumeSubHeadingListEnd
	
	
	%-----------PROJECTS-----------
	\section{Проекты}
	%Sample
	% \resumeSubHeadingListStart
	%   \resumeProjectHeading
	%       {\textbf{Gitlytics} $|$ \emph{Python, Flask, React, PostgreSQL, Docker}}{June 2020 -- Present}
	%       \resumeItemListStart
	%         \resumeItem{Developed a full-stack web application using with Flask serving a REST API with React as the frontend}
	%         \resumeItem{Implemented GitHub OAuth to get data from user’s repositories}
	%         \resumeItem{Visualized GitHub data to show collaboration}
	%         \resumeItem{Used Celery and Redis for asynchronous tasks}
	%       \resumeItemListEnd
	
	\resumeSubHeadingListStart
	\resumeSubheading
	{\textbf{PRAssign}  $|$ 
		\normalfont{\emph{\ulhref{https://github.com/UsatovPavel/PRAssign}{Go REST API} + {\ulhref{https://github.com/UsatovPavel/AsyncFactorial}{Scala Kafka consumers}:}
			}
	}}{ноябрь 2025 - январь 2026}
	{Получение задач через Go, передача через Kafka, вычисление на Scala, отслеживание в БД} {Индивидуальный проект}
	%\resumeItemListStart
	%\resumeItem 
	%\resumeItem{}{API для управления PR с проверкой прав и статистикой}
	%\resumeItem{}{Асинхронный пайплайн обработки задач: Go API $\rightarrow$ Kafka $\rightarrow$ Scala consumer $\rightarrow$ Kafka$\rightarrow$ PostgreSQL;}
	%\resumeItem{}{Контейнеризированный запуск/тестирование и фоновая очистка}
	%\resumeItem{}{Масштабирование consumer-групп и API (Nginx), параллельная обработка на Cats Effect/FS2}
	%\resumeItemListEnd
	\par\noindent\textbf{Технологии}:  Scala, Go, 
	Kafka, Cats Effect 3, PostgreSQL, Nginx, ScalaTest, Docker, k6%,Go Gin golang-migrate
	
	\resumeSubheading
	{\ulhref{https://github.com/hse-project-Java-2025}{Приложение «TimeTamer»}}{февраль 2025 - август 2025} 
	{Java server$+$Kotlin Android приложение-календарь с элементами AI и геймификации} {Командный проект}
	\resumeItemListStart
	\item {Общие задания, статистика, достижения, push-уведомления, аутентификация}
	\item {Взаимодействие с сервером через Retrofit и WorkManager}
	\item{Интеграция AI-ассистента (ChatGPT + Whisper) для предложений по задачам и голосового ввода.}
	\resumeItemListEnd
	\textbf{Технологии}: Java, Spring Boot, Kotlin, Jetpack Compose, Retrofit, PostgreSQL, Docker, JUnit 4/5,  Mockito
	
	\resumeSubheading
	{\ulhref{https://github.com/UsatovPavel/Voevoda}{Игра «Voevoda»}}{январь 2024 - июнь 2024}
	{2d topdown RTS: управление генералом, захват городов, наём армии и сражения с ИИ}{Командный проект}
	\resumeItemListStart
	\resumeItem{}{Генерация городов и оппонентов, поведение врагов, виджеты армии ИИ-оппонентов;  туман войны и сражения}
	\resumeItemListEnd
	\textbf{Технологии}: C++, Unreal Engine 4
	
	\resumeSubheading
	{\ulhref{https://github.com/FUYOH666/scanovich-webUI/}{Хакатон MTS True Tech}}
	{2026}
	{разработка оркестратора нейросетей на Python, интеграция с WebUI}{Командный проект}
	
	\resumeSubheading
	{Учебные задания}{}{}{}% ТУПО: либо антистек, либо 
	\resumeItemListStart
	% \resumeItem{C++}{библиотека для тестов с использованием макросов}
	\resumeItem{Java}{реализация Cli аналога Git}
	\resumeItem{Scala}{\ulhref{https://github.com/UsatovPavel/ZIO-Notification-Service}{ZIO REST API} для уведомлений с учётом часовых поясов}% с тестированием через TestClient
	\resumeItem{Python}{\ulhref{https://github.com/UsatovPavel/Usatov-FL-HSE/tree/task4-dev}{ANTLR-парсер} регулярных выражений}%\ulhref{https://github.com/UsatovPavel/Pop_ML_Forecast10}{ML: Предсказание населения стран на 10-летнем горизонте}
	\resumeItemListEnd
	\resumeSubHeadingListEnd
	
	%
	%-----------PROGRAMMING SKILLS-----------
	\section{Навыки}
	\textbf{Языки}{: Java, Scala, Kotlin} \\ %, Python, C++)} \\
	\textbf{Frameworks}{: Spring Boot, JUnit, Mockito, Cats Effect} \\
	\textbf{Технологии и инструменты}{: Git, Docker, PostgreSQL, Kafka, Kubernetes, gRPC, Grafana, Prometheus, Gradle} \\%Nginx, k6, p2p, ZIO, FS
	\textbf{Навыки}{: тестирование (интеграционное/e2e/нагрузочное/UI), алгоритмы и структуры данных, работа по Kanban}%проектирование архитектур(MVVM/клиент-сервер/микросервисная)}%<- ТУПО
% Jetpack Compose, Retrofit, WorkManager <- если на Mobile/fullstack добавить
% Unreal Engine <- если на Gamedev


%\section{Дополнительная деятельность}
%\resumeSubHeadingListStart
%\resumeSubheading
%{Кировская летняя метапредметная школа}{Вишкиль}
%{Педагог доп. образования (математика)}{Июль 2024, Июль 2025}
%\resumeSubHeadingListEnd

%-------------------------------------------
\end{document}